

\def\norm#1{\left\|#1\right\|}

\newcommand{\D}{D}
\newcommand{\Lun}{L}
\newcommand{\Lsym}{L_{\text{sym}}}
\newcommand{\Lrw}{L_{\text{rw}}}
\renewcommand{\bar}{\overline} %

\DeclareMathOperator{\Tr}{Tr}%
\DeclareMathOperator{\vol}{vol} %
\DeclareMathOperator{\Ratiocut}{RatioCut}
\DeclareMathOperator{\Ncut}{Ncut}
\DeclareMathOperator{\Cut}{cut}
\DeclareMathOperator{\Mincut}{MinCut}
\DeclareMathOperator{\MinmaxCut}{MinMaxCut}

 

%\usepackage{xcolor}% http://ctan.org/pkg/xcolor
%\usepackage{colortbl}% http://ctan.org/pkg/colortbl
%%\def\pedro#1{\textcolor{magenta}{#1}} 
%\usepackage[utf8]{inputenc}
%\usepackage{amsfonts}
%%%%%%%%%%%%%%%%%%%%%%%%%%%%%%%%%%%%%%%%%%%
%%%%%%%%%%%%%%%%%%%%%%%%%%%%%%%%%%%%%%%%%%%
%%%%%%%%%%%%%%%%%%%%%%%%%%%%%%%%%%%%%%%%%%%
%   
%\pagestyle{plain}
%\usepackage[numbers]{natbib}
% Comment this out before submission to drop the page numbers 

%%%%
%%%%
%%%  From Dropbox
%%%


\newcommand\x{\times}
\newcommand\bigzero{\makebox(0,0){\text{\huge0}}}
\newcommand*{\bord}{\multicolumn{1}{c|}{}}
 
\newcommand{\Ln}{L^{-}}
\newcommand{\Lp}{L^{+}}
\newcommand{\Dn}{D^{-}}
\newcommand{\Dp}{D^{+}}
 
\newcommand{\Tbar }{ \bar{T} }
 
\newcommand{\ELn }{   \expec{[\Ln]}  }
\newcommand{\ELp }{   \expec{[\Lp]}  }

\newcommand{\EDp }{   \expec{[\Dp]}  }
\newcommand{\EDn }{   \expec{[\Dn]}  }

    
\newcommand{\An}{A^{-}}
\newcommand{\Ap}{A^{+}}
\newcommand{\EAp }{   \expec{[\Ap]}  }
\newcommand{\EAn }{   \expec{[\An]}  }



\usepackage{tikz}
\usetikzlibrary{matrix,decorations.pathreplacing,calc}

% Set various styles for the matrices and braces. It might pay off to fiddle around with the values a little bit
\pgfkeys{tikz/mymatrixenv/.style={decoration=brace,every left delimiter/.style={xshift=3pt},every right delimiter/.style={xshift=-3pt}}}
\pgfkeys{tikz/mymatrix/.style={matrix of math nodes,left delimiter=[,right delimiter={]},inner sep=2pt,column sep=1em,row sep=0.5em,nodes={inner sep=0pt}}}
\pgfkeys{tikz/mymatrixbrace/.style={decorate,thick}}
\newcommand\mymatrixbraceoffseth{0.5em}
\newcommand\mymatrixbraceoffsetv{0.2em}

% Now the commands to produce the braces. (I'll explain below how to use them.)
\newcommand*\mymatrixbraceright[4][m]{
    \draw[mymatrixbrace] ($(#1.north west)!(#1-#3-1.south west)!(#1.south west)-(\mymatrixbraceoffseth,0)$)
        -- node[left=2pt] {#4} 
        ($(#1.north west)!(#1-#2-1.north west)!(#1.south west)-(\mymatrixbraceoffseth,0)$);
}
\newcommand*\mymatrixbraceleft[4][m]{
    \draw[mymatrixbrace] ($(#1.north east)!(#1-#2-1.north east)!(#1.south east)+(\mymatrixbraceoffseth,0)$)
        -- node[right=2pt] {#4} 
        ($(#1.north east)!(#1-#3-1.south east)!(#1.south east)+(\mymatrixbraceoffseth,0)$);
}
\newcommand*\mymatrixbracetop[4][m]{
    \draw[mymatrixbrace] ($(#1.north west)!(#1-1-#2.north west)!(#1.north east)+(0,\mymatrixbraceoffsetv)$)
        -- node[above=2pt] {#4} 
        ($(#1.north west)!(#1-1-#3.north east)!(#1.north east)+(0,\mymatrixbraceoffsetv)$);
}
\newcommand*\mymatrixbracebottom[4][m]{
    \draw[mymatrixbrace] ($(#1.south west)!(#1-1-#3.south east)!(#1.south east)-(0,\mymatrixbraceoffsetv)$)
        -- node[below=2pt] {#4} 
        ($(#1.south west)!(#1-1-#2.south west)!(#1.south east)-(0,\mymatrixbraceoffsetv)$);
}

  
\newcommand{\mydef}{\stackrel{\text{def}}{=}}
\newcommand{\mydeff}{ \coloneqq }
 
\newcommand{\omo}{ \mb{1}_{C, \bar{C}}  }

% \usepackage{nicematrix}
% \NiceMatrixOptions{transparent,nullify-dots}
 
\newcommand{\tce}{   \tilde{c}_{\varepsilon}  }

%-----------------------------------------
% Some commands defined by Hemant
%----------------------------------------- 
\newcommand{\EA }{   \expec{[A]}  }
% ????  \newcommand{\Dbar }{ \bar{D} } 
\newcommand{\EDbar }{   \expec{[\Dbar]}  }
\newcommand{\Dbarn }{ \bar{D}^{-} }
\newcommand{\Dbarp }{ \bar{D}^{+} }
\newcommand{\Lbar }{ \bar{L} }
\newcommand{\ELbar }{   \expec{[\Lbar]}  }

% \newcommand{\mb}[1]{\mbox{\boldmath$#1$}}
\newcommand{\ones }{ \mb{1} }

\newcommand{\infovec }{w}
\newcommand{\pert }{\tri}

\newcommand{\Gn }{ G^{-} }
\newcommand{\Gp }{ G^{+} }

\newcommand{\taup }{ \tau^{+} }
\newcommand{\taun }{ \tau^{-} }



% Miscellaneous symbols 
\newcommand{\tri}{\ensuremath{\bigtriangleup}}

\newcommand{\ceil}[1]{\lceil{#1}\rceil}
\newcommand{\floor}[1]{\lfloor{#1}\rfloor}
\newcommand{\set}[1]{\left\{{#1}\right\}}
\newcommand{\dotprod}[2]{\langle#1,#2\rangle}
\newcommand{\est}[1]{\widehat{#1}}
\newcommand{\expec}{\ensuremath{\mathbb{E}}}
\newcommand{\grad}{\ensuremath{\nabla}}
\newcommand{\hess}{\ensuremath{\nabla^2}}
\newcommand{\matR}{\ensuremath{\mathbb{R}}}
\newcommand{\matZ}{\ensuremath{\mathbb{Z}}}
\newcommand{\argmax}[1]{\underset{#1}{\operatorname{argmax}}}
\newcommand{\argmin}[1]{\underset{#1}{\operatorname{argmin}}}
\newcommand{\ubar}[1]{\underaccent{\bar}{#1}}
\newcommand{\prob}{\ensuremath{\mathbb{P}}}
\newcommand{\Linfnorm}{L_{\infty}}
\newcommand{\vp}{\ensuremath{v^{\prime}}}
\newcommand{\vpp}{\ensuremath{v^{\prime\prime}}}
\newcommand{\tp}{\ensuremath{t^{\prime}}}
\newcommand{\lp}{\ensuremath{l^{\prime}}}
\newcommand{\qp}{\ensuremath{q^{\prime}}}


\newcommand{\red}[1]{\textcolor{red}{#1}}
\newcommand{\blue}[1]{\textcolor{blue}{#1}}
